% Chapter 3 -- Configuring Your System (the CONFIG.SYS commands; IBM 3.10
% Reference Ch. 4 "Configuring Your System"; AN-01 "Configuration Command"
% entries).  The entry list is the shipped directive table of
% os/milton/config_sys.c config_sys_parse: BREAK BUFFERS DEVICE FILES INSTALL
% LASTDRIVE SHELL.  What each one does is os/milton/sysinit.c (FILES applied
% to the SFT cap of 20; BREAK applied; DEVICE=ANSI.SYS enables the ANSI
% console, int21_set_ansi_gate; the others recorded "accepted(deferred)").
% Standard settings: spec/dos_config_sys_baseline.txt (LOCKED).  The file is
% read from A:\CONFIG.SYS, first 1024 bytes, to the last whole line.
% GATES: test-config-sys test-config-fuzz test-sysinit test-sysinit-oversize test-ansi
\Gates{test-config-sys, test-config-fuzz, test-sysinit, test-sysinit-oversize, test-ansi}
\chapter{Configuring Your System}

\section{The CONFIG.SYS File}
\index{CONFIG.SYS@\texttt{CONFIG.SYS}}
When InitechDOS starts, before it displays the command prompt, it reads the
file \file{CONFIG.SYS} in the root directory of drive~A.  The file holds
\emph{configuration commands}, one to a line, that set up the system: how
many files can be open at once, which device drivers are used, and so on.

Each line has the form
\SyntaxLine{\ph{command}\kw{=}\ph{value}}
in capital or small letters.  Spaces around the equal sign are ignored.
Blank lines, lines that begin with a semicolon (\texttt{;}), and lines
InitechDOS does not recognize are ignored, so you can use semicolon lines
for remarks.  A value outside the allowed range is changed to the nearest
allowed value.

InitechDOS reads the first 1,024 characters of the file, up to the end of
the last complete line.  Changes to \file{CONFIG.SYS} take effect the next
time you start the system.

\section{Creating CONFIG.SYS}
\file{CONFIG.SYS} is an ordinary text file.  You can create it from the
keyboard with \cmd{COPY CON}:
\begin{TypeIn}
COPY CON CONFIG.SYS
FILES=20
BREAK=OFF
\end{TypeIn}
and then press \key{Ctrl-Z} and \key{Enter}.  Restart the system to use the
new file.

\section{The Standard Settings}
\index{standard settings}
If drive~A has no \file{CONFIG.SYS} file, InitechDOS uses these standard
settings, exactly as if they were in the file:
\begin{TypeIn}
FILES=20
BUFFERS=20
LASTDRIVE=Z
DEVICE=ANSI.SYS
DEVICE=INITNET.SYS
INSTALL=SHARE.EXE
SHELL=COMMAND.COM /P /E:512
\end{TypeIn}
The CTRL-BREAK check is on unless a \file{CONFIG.SYS} file says
\texttt{BREAK=OFF}.

The entries that follow describe each configuration command.

% BREAK= -- config_sys.c (ON/OFF case-insensitive; other values ignored);
% sysinit.c int21_set_break_flag (default ON).
\begin{DosEntry}{BREAK}{Control-Break}{config}
\Purpose Sets the CTRL-BREAK check when the system starts.
\Format \SyntaxLine{\kw{BREAK=ON}\texttt{|}\kw{OFF}}
\Remarks
\kw{ON} turns the check on and \kw{OFF} turns it off.  Without a
\texttt{BREAK} line the check is on.  Any other value is ignored.

The setting can be changed and displayed at any time with the internal
command \cmd{BREAK} (Chapter~2), which describes what the check does.
\Example
\begin{TypeIn}
BREAK=OFF
\end{TypeIn}
\end{DosEntry}

% BUFFERS= -- config_sys.c (1..99, clamped); sysinit.c records it
% "accepted(deferred)": there is no buffer pool in this version.
\begin{DosEntry}{BUFFERS}{}{config}
\Purpose Sets the number of disk buffers.
\Format \SyntaxLine{\kw{BUFFERS=}\ph{n}}
\Remarks
\Param{n} is a number from 1 to 99.  The standard setting is 20.

InitechDOS checks the value and records it when the system starts.  This
version reads and writes disk sectors directly, without a pool of disk
buffers, so the number does not change the way the system runs.
\Example
\begin{TypeIn}
BUFFERS=20
\end{TypeIn}
\end{DosEntry}

% DEVICE= -- config_sys.c (up to 8 lines; first token kept); sysinit.c
% (ANSI.SYS -> ansi-enabled; any other driver accepted(deferred)); the ANSI
% console: int21.c (ESC[ sequences, gate test-ansi).
\begin{DosEntry}{DEVICE}{}{config}
\Purpose Names a device driver to be installed when the system starts.
\Format \SyntaxLine{\kw{DEVICE=}\opt{\ph{d}\kw{:}}\opt{\ph{path}}\ph{filename}\opt{\kw{.}\ph{ext}}
  \opt{\ph{parameters}}}
\Remarks
\Param{filename} names the driver.  Up to eight \texttt{DEVICE} lines are
used.

\texttt{DEVICE=ANSI.SYS} installs the extended screen and keyboard control
of the console: a program can then move the cursor, clear the screen and
set display attributes by writing escape sequences to the display.  The
ANSI.SYS support is part of InitechDOS; no file of that name is needed on
the disk.  It is installed by the standard settings.

This version of InitechDOS supports ANSI.SYS only.  Any other
\texttt{DEVICE} line is checked and recorded when the system starts, and the
driver it names is not loaded.  The console, printer, serial port, clock and
\texttt{NUL} devices are always present and need no \texttt{DEVICE} line.
\Example
\begin{TypeIn}
DEVICE=ANSI.SYS
\end{TypeIn}
\end{DosEntry}

% FILES= -- config_sys.c (8..255 clamped); sysinit.c sft_set_files_limit
% (table of 20: cap=20).
\begin{DosEntry}{FILES}{}{config}
\Purpose Sets the number of files that can be open at one time.
\Format \SyntaxLine{\kw{FILES=}\ph{n}}
\Remarks
\Param{n} is a number from 8 to 255.  The standard setting is 20.

InitechDOS keeps a table of 20 open files for the whole system, so a value
larger than 20 is treated as 20.
\Example
\begin{TypeIn}
FILES=20
\end{TypeIn}
\end{DosEntry}

% INSTALL= -- config_sys.c (up to 8 lines; first token); sysinit.c
% "accepted(deferred)".
\begin{DosEntry}{INSTALL}{}{config}
\Purpose Names a program to be loaded when the system starts.
\Format \SyntaxLine{\kw{INSTALL=}\opt{\ph{d}\kw{:}}\opt{\ph{path}}\ph{filename}\opt{\kw{.}\ph{ext}}
  \opt{\ph{parameters}}}
\Remarks
\Param{filename} names the program.  Up to eight \texttt{INSTALL} lines are
used.

This version of InitechDOS checks and records each \texttt{INSTALL} line
when the system starts; it does not load the program.  The standard
settings name \file{SHARE.EXE}.
\Example
\begin{TypeIn}
INSTALL=SHARE.EXE
\end{TypeIn}
\end{DosEntry}

% LASTDRIVE= -- config_sys.c (A..Z); sysinit.c "accepted(deferred)".
\begin{DosEntry}{LASTDRIVE}{}{config}
\Purpose Sets the last drive letter InitechDOS can use.
\Format \SyntaxLine{\kw{LASTDRIVE=}\ph{x}}
\Remarks
\Param{x} is a letter from A to Z.  The standard setting is Z.

InitechDOS checks and records the letter when the system starts.  In this
version the disk is always drive~A, so the setting does not change the
drives you can use.
\Example
\begin{TypeIn}
LASTDRIVE=E
\end{TypeIn}
\end{DosEntry}

% SHELL= -- config_sys.c (whole tail kept, 64 characters); sysinit.c
% records it; COMMAND.COM is resident (command.c command_repl, COMSPEC
% A:\COMMAND.COM); environment fixed at 512 bytes (env.h ENV_ARENA_MAX).
\begin{DosEntry}{SHELL}{}{config}
\Purpose Names the command processor and its parameters.
\Format \SyntaxLine{\kw{SHELL=}\opt{\ph{d}\kw{:}}\opt{\ph{path}}\ph{filename}\opt{\kw{.}\ph{ext}}
  \opt{\ph{parameters}}}
\Remarks
The standard setting is \texttt{SHELL=COMMAND.COM /P /E:512}: the command
processor \file{COMMAND.COM}, made permanent (\kw{/P}), with an environment
of 512 characters (\kw{/E:512}).

InitechDOS checks and records the line when the system starts.  In this
version the command processor is always \file{COMMAND.COM}, it is always
permanent, and its environment is always 512 characters, whatever the
\texttt{SHELL} line says.
\Example
\begin{TypeIn}
SHELL=COMMAND.COM /P /E:512
\end{TypeIn}
\end{DosEntry}
